\section{Charakteristiky polohy}\label{sec:ai-charakteristiky-polohy}

Charakteristiky polohy nahrazují celý statistický soubor jedním číslem, které popisuje jeho typickou nebo střední hodnotu. Nejčastěji budeme používat aritmetický průměr a~medián. Každá z~těchto charakteristik zdůrazňuje jinou vlastnost dat.

\subsection{Aritmetický a~vážený průměr}

\begin{definition}\label{def:ai-prumer}
    \emph{Aritmetický průměr} statistického souboru
    \(\mathbf{X}=(x_1,\ldots,x_n)\) je číslo
    \[
        \average{\mathbf{X}}
        =
        \frac{1}{n}\sum_{i=1}^{n}x_i.
    \]
\end{definition}

Průměr zohledňuje všechny hodnoty a~každé přisuzuje stejný význam. Nemusí být jednou z~naměřených hodnot. Je také citlivý na odlehlá pozorování: jediná mimořádně vysoká nebo nízká hodnota jej může výrazně posunout.

\begin{proposition}\label{prop:ai-prumer-vlastnosti}
    Pro statistické soubory \(\mathbf{X},\mathbf{Y}\) délky \(n\) a~čísla \(a,b\in\R\) platí
    \[
        \sum_{i=1}^{n}(x_i-\average{\mathbf{X}})=0
        \qquad\text{a}\qquad
        \average{a\mathbf{X}+b\mathbf{Y}}
        =
        a\average{\mathbf{X}}+b\average{\mathbf{Y}}.
    \]
    Aritmetický průměr navíc minimalizuje součet čtvercových odchylek:
    \[
        \average{\mathbf{X}}
        =
        \argmin_{c\in\R}
        \sum_{i=1}^{n}(x_i-c)^2.
    \]
\end{proposition}

\begin{proof}
    První dvě rovnosti získáme přímým dosazením definice:
    \[
        \sum_{i=1}^{n}(x_i-\average{\mathbf{X}})
        =
        \sum_{i=1}^{n}x_i-n\average{\mathbf{X}}
        =
        0
    \]
    a
    \[
        \average{a\mathbf{X}+b\mathbf{Y}}
        =
        \frac1n\sum_{i=1}^{n}(ax_i+by_i)
        =
        a\average{\mathbf{X}}+b\average{\mathbf{Y}}.
    \]
    Pro libovolné \(c\in\R\) dále platí
    \begin{align*}
        \sum_{i=1}^{n}(x_i-c)^2
        &=
        \sum_{i=1}^{n}
        \bigl((x_i-\average{\mathbf{X}})
        +(\average{\mathbf{X}}-c)\bigr)^2\\
        &=
        \sum_{i=1}^{n}(x_i-\average{\mathbf{X}})^2
        +2(\average{\mathbf{X}}-c)
        \underbrace{\sum_{i=1}^{n}(x_i-\average{\mathbf{X}})}_{0}
        +n(\average{\mathbf{X}}-c)^2.
    \end{align*}
    První člen nezávisí na \(c\) a~poslední člen je nejmenší právě pro
    \(c=\average{\mathbf{X}}\).
\end{proof}

Ne vždy mají všechny hodnoty stejný význam. Máme-li vektor kladných vah
\(\mathbf{w}=(w_1,\ldots,w_n)\), definujeme vážený průměr vztahem
\[
    \average{\mathbf{X}}_{\mathbf{w}}
    =
    \frac{\sum_{i=1}^{n}w_ix_i}{\sum_{i=1}^{n}w_i}.
\]
Po zavedení normalizovaných vah
\[
    \omega_i=\frac{w_i}{\sum_{j=1}^{n}w_j}
\]
platí \(\omega_i>0\), \(\sum_{i=1}^{n}\omega_i=1\) a
\[
    \average{\mathbf{X}}_{\mathbf{w}}
    =
    \sum_{i=1}^{n}\omega_ix_i.
\]
Vážený průměr proto vždy leží mezi nejmenší a~největší hodnotou souboru.

\subsection{Medián}

Seřazené hodnoty souboru budeme označovat
\[
    x_{(1)}\leqslant x_{(2)}\leqslant\cdots\leqslant x_{(n)}.
\]

\begin{definition}\label{def:ai-median}
    \emph{Medián} statistického souboru \(\mathbf{X}\) je číslo
    \[
        \median(\mathbf{X})
        =
        \begin{cases}
            x_{((n+1)/2)}, & \text{je-li \(n\) liché},\\[1mm]
            \dfrac{x_{(n/2)}+x_{(n/2+1)}}{2},
            & \text{je-li \(n\) sudé}.
        \end{cases}
    \]
\end{definition}

Medián dělí uspořádaná data na dolní a~horní polovinu. Na rozdíl od průměru jej jednotlivá odlehlá hodnota obvykle ovlivní jen málo. Pro soubor
\[
    \mathbf{X}=(2,3,4,5,100)
\]
například dostáváme
\[
    \median(\mathbf{X})=4,
    \qquad
    \average{\mathbf{X}}=22{,}8.
\]

\begin{figure}[ht]
    \centering
    \begin{tikzpicture}[x=1.05cm,y=1cm,every node/.style={font=\normalsize}]
    \draw[diredge] (0,0) -- (10.4,0) node[right] {\(x\)};
    \foreach \x in {2,3,4,5,10}{
        \fill[blue!75!black] (\x,0) circle (2.4pt);
    }
    \node[above] at (2,0.12) {\(2\)};
    \node[above] at (3,0.12) {\(3\)};
    \node[above] at (4,0.12) {\(4\)};
    \node[above] at (5,0.12) {\(5\)};
    \node[above] at (10,0.12) {\(100\)};
    \draw[green!60!black,very thick] (4,-0.45) -- (4,0.45);
    \node[green!50!black,below] at (4,-0.48) {\(\median(\mathbf{X})\)};
    \draw[red!75!black,very thick] (6.1,-0.45) -- (6.1,0.45);
    \node[red!75!black,below] at (6.1,-0.48) {\(\average{\mathbf{X}}\)};
    \draw[decorate,decoration={brace,amplitude=5pt},thick]
        (5.15,0.75) -- (9.9,0.75)
        node[midway,above=5pt]{odlehlá hodnota posouvá průměr};
\end{tikzpicture}

    \caption{Odlehlá hodnota posune průměr podstatně více než medián.}
    \label{fig:ai-prumer-median}
\end{figure}

\tipbox{Průměr je přirozený pro algebraické výpočty a~pro metodu nejmenších čtverců. Medián je odolnější vůči odlehlým hodnotám. Volba charakteristiky proto závisí na povaze dat a~na otázce, kterou chceme zodpovědět.}
